\editionPage[\pdfbookmark[0]{Comprendre pourquoi}{bm-preuves-debut-parent}]{preuves-debut}{Comprendre pourquoi}{Pourquoi on compte les décimales d’un produit}{Les explications ci-dessous éclairent des règles utilisées dans le livre.}
\begin{blocEdition}{preuves-debut}
\editionEtape{Écrire avec des dixièmes}$2{,}3$, c’est 23 dixièmes : $2{,}3=23\times0{,}1$. De même, $1{,}4=14\times0{,}1$.\editionEtape{Changer l’ordre}On peut changer l’ordre des facteurs d’un produit et les regrouper : $2{,}3\times1{,}4=23\times14\times0{,}1\times0{,}1$.\editionEtape{Un dixième de dixième}$0{,}1\times0{,}1$, c’est un dixième d’un dixième, soit un centième : $0{,}01$. Donc $2{,}3\times1{,}4=322\times0{,}01=3{,}22$.\label{fig-preuves-debut-preuve-produit}\figureGeo{preuve-produit}
\editionRetenir{Un facteur à une décimale apporte des dixièmes, à deux décimales des centièmes : le nombre de décimales du produit est la somme des nombres de décimales des facteurs.}
\end{blocEdition}
\editionPage{preuves-2}{Comprendre pourquoi}{Pourquoi un tiers n’est pas un nombre décimal}{Les explications ci-dessous éclairent des règles utilisées dans le livre.}
\begin{blocEdition}{preuves-2}
\editionEtape{Supposer le contraire}Si $\frac13$ était décimal, il s’écrirait comme une fraction décimale, par exemple $\frac{a}{1\,000}$ : on aurait alors $3\times a=1\,000$.\editionEtape{Une impossibilité}Or $1\,000$ n’est pas dans la table de 3 : la somme de ses chiffres vaut 1. De même, $10$, $100$, $10\,000$… ne sont jamais des multiples de 3.\editionEtape{Conclure}Aucune fraction décimale n’est égale à $\frac13$ : ce nombre n’est pas décimal. Sa division ne s’arrête jamais : $0{,}333\ldots$ On utilise des valeurs approchées, comme $0{,}33$.\editionRetenir{Une fraction peut représenter un nombre qui n’a pas d’écriture décimale exacte. Il est alors exact de l’écrire en fraction, et approché de l’écrire avec une virgule.}
\end{blocEdition}
\editionPage{preuves-3}{Comprendre pourquoi}{Pourquoi la somme des angles d’un triangle vaut 180°}{Les explications ci-dessous éclairent des règles utilisées dans le livre.}
\begin{blocEdition}{preuves-3}
\editionEtape{Une expérience}On découpe un triangle en papier, puis on déchire ses trois coins. Placés côte à côte, sommets réunis, les trois angles forment un angle plat.\label{fig-preuves-3-preuve-triangle}\figureGeo{preuve-triangle}
\editionEtape{Avec n’importe quel triangle}L’expérience réussit avec un triangle long et fin comme avec un triangle équilatéral : la somme vaut toujours $180^\circ$. Pour un triangle équilatéral, on retrouve $3\times60^\circ=180^\circ$.\editionEtape{Vers une preuve}Au collège, on démontrera ce résultat pour tous les triangles à la fois, avec des propriétés des droites parallèles : une expérience convainc, une démonstration prouve.\editionRetenir{Un triangle ne peut pas avoir deux angles obtus, ni deux angles droits : leur somme dépasserait déjà $180^\circ$.}
\end{blocEdition}
\editionPage{preuves-4}{Comprendre pourquoi}{Pourquoi les trois médiatrices se coupent}{Les explications ci-dessous éclairent des règles utilisées dans le livre.}
\begin{blocEdition}{preuves-4}
\editionEtape{Deux médiatrices}Dans le triangle ABC, les médiatrices de [AB] et de [BC] se coupent en un point O.\editionEtape{Des distances égales}O est sur la médiatrice de [AB], donc $OA=OB$. O est sur la médiatrice de [BC], donc $OB=OC$. Par conséquent $OA=OC$.\editionEtape{La troisième}Puisque $OA=OC$, le point O est à égale distance de A et de C : il est sur la médiatrice de [AC], par la propriété réciproque. Les trois médiatrices passent par O.\editionRetenir{Comme $OA=OB=OC$, le cercle de centre O et de rayon OA passe par les trois sommets : c’est le cercle circonscrit au triangle.}
\end{blocEdition}
